\ficha{Topik (1987), The Political Economy of the Brazilian State}{topik1987}

\begin{identificacao}
TOPIK, Steven. \emph{The Political Economy of the Brazilian State, 1889-1930}.
Austin: University of Texas Press, 1987. (Latin American Monographs, Institute
of Latin American Studies, The University of Texas at Austin, n. 71). ISBN
0-292-76500-2.

Livro de história econômica, em inglês, 168 páginas de texto e notas a partir da
p. 169. Leitura seletiva: capítulos 1 (p. 1-26), 2 (\emph{The Financial System},
p. 27-58) e 3 (\emph{The Defense of Coffee}, p. 59-92) por inteiro, capítulo 6
(p. 161-168) por alto; capítulos 4 e 5 não fichados. Exemplar consultado em PDF
da University of Texas Press, série LLILAS Latin American Monographs, com a
paginação da primeira edição de 1987.
\end{identificacao}

\bloco{Problema e tese}
A pergunta é quem exerceu o poder de Estado na Primeira República e por que um
regime que se dizia liberal terminou com presença econômica tão grande. A
resposta é a autonomia relativa: o Estado não foi correia de transmissão dos
fazendeiros de café nem casta burocrática patrimonial, e a margem de manobra que
teve veio menos do conflito entre classes internas do que do peso dos
investidores estrangeiros (p. 164). A intervenção foi resultado do choque entre
frações da classe dominante e credores externos, e é a estrutura da economia, não
a intenção dos atores, que a explica (p. 26).

\bloco{Argumento}
O livro inteiro sustenta que o intervencionismo estatal brasileiro não nasce em
1930 nem da falência da economia exportadora, mas é produto dessa mesma economia
exportadora. Topik mede a afirmação em quatro setores, finanças, café, ferrovias
e indústria, e conclui que Vargas herdou instituições públicas robustas, não o
esqueleto frágil da historiografia corrente (p. 161). O paradoxo que fecha o
livro: o Estado interveio para salvar o modelo laissez faire e, ao intervir, o
destruiu (p. 168).

\begin{enumerate}[leftmargin=*,nosep]
\item A autonomia relativa é definida contra dois adversários. Contra a
dependência, que faz do Estado liberal um ``messenger boy'' da classe dominante
(p. 3). Contra o patrimonialismo: os administradores são antes comitê de gestão
dos negócios comuns da burguesia, na fórmula de Marx que Topik adota, e rompem
com os fazendeiros porque respondem pelo longo prazo do sistema (p. 163).

\chave{The state's relative autonomy from the will of fazendeiros is explained by
the major role foreign investors played in the Brazilian economy more than by
contending class interests within Brazil.}{p. 164}

\item Quase todos os ministros da Fazenda quiseram, na frase de Furtado que Topik
adota, submeter o sistema econômico às regras monetárias da Europa, e nenhum
entregou: o mil-réis caiu de 26 pence em 1889 para 5 pence no fim da República e
o padrão-ouro nunca se impôs (pp. 27, 50). Ao mesmo tempo o Estado passou a
controlar câmbio, crédito comercial e hipoteca por instituições públicas que
minaram os bancos estrangeiros (p. 27).

\item O enforcement externo tem nome e data. Quando Campos Sales, presidente
eleito, sondou uma moratória, os Rothschild responderam que, além da perda
completa do crédito, a medida poderia afetar a soberania do país e chegar ao
extremo da invasão estrangeira (p. 36). O funding loan de 1898 consolidou a
dívida, obrigou a queimar papel-moeda, hipotecou a alfândega do Rio e pôs bancos
estrangeiros a fiscalizar a execução; comerciantes e investidores estrangeiros o
apoiaram, fazendeiros de café e industriais o combateram, e os consumidores
urbanos se dividiram entre a objeção ao controle britânico e a expectativa de
câmbio melhor (p. 37).

\item O Banco do Brasil de 1905 é intervenção sem intenção nacionalista. Bulhões
buscou primeiro capital estrangeiro e desistiu porque os estrangeiros eram ``very
demanding'' (p. 39); os acionistas privados aceitaram porque esperavam
instituição capaz de exercer funções de banco central (p. 40). O banco virou
agente cambial do Tesouro, ``the supreme arbiter of the exchange market in
Brazil'', e ganhou o monopólio dos vales ouro que os importadores compravam para
pagar parte dos direitos alfandegários, o que lhe garantiu moeda estrangeira
contra os bancos estrangeiros; sua política era ditada pelo ministro da Fazenda,
não pelos acionistas (p. 42).

\item A Caixa de Conversão é o fundamento da reforma financeira de Afonso Pena:
casa de câmbio que emitia notas conversíveis contra moeda estrangeira lastreada
em ouro, a taxa melhor que a de mercado, e empregava as libras acumuladas para
manter o mil-réis mais de 50\% abaixo do par legal (p. 40). Rodrigues Alves a
vetara porque taxa abaixo do par afugentaria o investidor estrangeiro,
encareceria importações e taxaria a riqueza em forma monetária (p. 41). A taxa
ficou em torno de 15 pence; com o fechamento momentâneo de 1910 o câmbio saltou
acima de 18; os paulistas queriam voltar a 15 e o ministro queria 18,25,
preocupado com o fim da carência do funding em 1911; o acordo saiu em 16 pence e
durou até a corrida ao ouro de 1914 (p. 41). Topik nega que fosse medida radical,
já que só operava por mercado, e ao mesmo tempo a trata como via de volta ao
ouro: em 1912, 40\% do meio circulante era conversível, e o Brasil nunca mais
chegaria tão perto da conversibilidade completa (pp. 41-42).

\chave{paralyzed the work of his [Pena's] predecessors}{p. 41}

\chave{It is essential to guarantee the efficiency of the redemption policy [of
paper money] that the government intervene in the gold market to regularize the
functions of supply and demand}{p. 42}

\item A valorização amarra câmbio e café. O Convênio de Taubaté pedia empréstimo
de quinze milhões de libras com garantia federal, imposto proibitivo sobre novos
cafezais, bolsa pública e propaganda no exterior; a cláusula mais controversa
juntava a sobretaxa de três francos por saca à Caixa de Conversão, ambas
desenhadas para facilitar o pagamento do empréstimo (p. 67). Rodrigues Alves
atacou a proposta por causa da Caixa e porque acreditava, como os Rothschild, que
o país não tinha recursos para controlar o preço mundial (p. 68). A oposição dos
Rothschild não era conspiração: queriam papel ainda mais agressivo do governo
central, com desapropriação de terras e monopólio federal de três anos sobre a
compra, a venda e a exportação do café (p. 69).

\item Os financiadores concorrentes entram pela recusa dos Rothschild. São Paulo
começou a valorização em 1906 com um milhão de libras do Disconto Gesellschaft e,
achando os bancos europeus reticentes por causa daquela oposição, recorreu aos
próprios exportadores e importadores de café, com o consórcio de Sielcken e
compras pela casa Theodore Wille (p. 71). O empréstimo de quinze milhões de
libras de 1908 foi o primeiro grande empréstimo externo tomado fora dos
Rothschild em mais de cinquenta anos e marca a internacionalização dos títulos
brasileiros; em troca, os estoques passaram a um comitê de credores, com um só
representante paulista e arbitragem do Banco da Inglaterra (p. 72). Só depois da
Primeira Guerra os Rothschild perdem a posição de agente monopolista (p. 165).

\item O eixo do conflito é de fração, não de região. De um lado, credores e
investidores estrangeiros, o Ministério da Fazenda, importadores e consumidores
urbanos, por moeda forte e preços estáveis. Do outro, agricultores, industriais,
varejistas e alguns banqueiros, pela expansão. Os dois campos queriam fortalecer
a iniciativa privada e os dois produziram mais Estado (pp. 56-57). União contra
São Paulo é conflito medido: foi sob os presidentes paulistas de 1894 a 1906, que
valorizaram o mil-réis em 35\% e nacionalizaram o maior banco do país (p. 40),
que os apelos por ajuda ao café menos foram atendidos; depois, Bernardes vetou a
Defesa Permanente porque subsídio ao café era subsídio a São Paulo, ``a danger to
national integration'' (p. 77). Os paulistas podem ter recebido o que precisavam,
mas não o que queriam (p. 162), e a única concessão relevante do Congresso ao
Centro-Sul foi um mil-réis barato (p. 91).
\end{enumerate}

\bloco{Conceitos e definições operacionais}
Autonomia relativa: capacidade de contrariar a fração dominante em nome do longo
prazo do sistema exportador, sem deixar de representá-lo, medida pelos episódios
em que presidentes fazendeiros negaram ajuda ao café (pp. 66, 162, 164). Caixa de
Conversão: casa de câmbio emissora de notas conversíveis contra moeda estrangeira
lastreada em ouro, que fixa a taxa abaixo do par legal e opera apenas por compra
e venda no mercado (pp. 40-42); Topik não usa \emph{currency board} nem cita a
Caja argentina. Valorização: compra e retenção de estoques para sustentar preço,
viável pela natureza perene do café e pela participação brasileira de mais de
dois terços do mercado mundial (p. 68). Capitalismo de Estado: uso do Banco do
Brasil, instituição semiprivada, para captar depósitos e executar política
monetária sem despertar oposição à interferência estatal (p. 55).

\bloco{Evidência e método}
História econômica e política, sem econometria, teste ou contrafactual formal.
Usa relatórios ministeriais, balanços da Contadoria da República, o Anuário
Estatístico da DGE, o \emph{Retrospecto Commercial} do \emph{Jornal do
Commercio}, relatórios do Banco do Brasil e das ferrovias, arquivos pessoais e
despachos consulares norte-americanos. O autor avisa na abertura que os números
do período são frágeis e devem ser lidos como ordem de grandeza, com erro de
10\% ou até de um terço sem alterar suas conclusões (p. ix). O método é a
comparação entre setores, entre estados e com outros países latino-americanos,
mais a reconstrução do conflito político em torno de cada decisão; séries
próprias de gasto federal real sustentam a tese da centralização (pp. 20-21).

\bloco{Agentes e vertentes}
Ortodoxos: Rodrigues Alves, Prudente de Morais, Campos Sales, Murtinho e Bulhões,
escola manchesteriana que os adversários no Congresso chamavam de ``spencerian''
(p. 66). Intervencionistas por circunstância: Afonso Pena, Nilo Peçanha e Hermes
da Fonseca, os presidentes de 1906 a 1914 que, sem ligação direta com o complexo
cafeeiro, romperam com a prática ortodoxa e deram aos fazendeiros crédito e
mil-réis barato (p. 43). Rui Barbosa é o autor da reforma de 1890, cuja ruptura
com o padrão-ouro Topik atribui menos ao projeto desenvolvimentista do que à
necessidade de acomodar a elite imperial (p. 30). Serzedello Corrêa comparece
como coronel que dizia ser o Brasil ainda colônia em termos de interesse
econômico (p. 32) e como ministro que, apesar disso, defendeu orçamento
equilibrado e volta ao ouro (p. 33). Rui, Serzedello, Murtinho, Bulhões e Amaro
Cavalcanti são listados como as maiores autoridades em economia do país (p. 24).
Banqueiros: os Rothschild como agente monopolista até a Primeira Guerra, o
Disconto Gesellschaft, Sielcken, Theodore Wille, Farquhar, a Brazilian Warrant
Company e o Banco da Inglaterra como árbitro; no Brasil, Mayrink e o Conde de
Figueiredo. Discute com Frank, Furtado, Faoro, Peláez e Suzigan, Holloway, Delfim
Netto, Love e Wallerstein.

\bloco{Críticas e limites}
A Caixa de Conversão ocupa três páginas e meia e é tratada por fora: não há
limite de emissão, fundo de garantia, fundo de resgate, agência de Londres nem
custódia do ouro de 1911, e a suspensão de 1914 comparece em uma oração
subordinada (p. 41). A atribuição da frase sobre a Caixa ter paralisado a obra
dos antecessores está confusa no original, que apresenta Rodrigues Alves como
``minister of finance'' quando ele acabara de deixar a presidência (p. 41); o
nome de Afonso Pena aparece grafado ``Alfonso'' (p. 40). Topik não discute o
debate público sobre o nível da taxa a não ser pelo episódio de 1910 e não trata
metalismo e papelismo como categorias, o que faz dele fonte de estrutura e de
conflito, não de vocabulário de época. A tese da autonomia relativa se apoia em
episódios, e o próprio autor admite que a aparente independência dos
administradores refletia o esforço de conciliar interesses contraditórios
(p. 56), o que a reduz a efeito de equilíbrio entre pressões.

\begin{dialogo}
\item Enforcement externo no discurso interno. Procurar se os Rothschild, o
funding de 1898 e a alfândega hipotecada são invocados como advertência contra a
Caixa ou como razão para adotá-la. Topik trata a ameaça de 1898 como enforcement
explícito, com menção a soberania e invasão (p. 36). Testa se o argumento do
credor externo circula no debate público ou fica restrito à negociação.

\item Concorrência entre financiadores. Procurar menções ao Disconto
Gesellschaft, a Sielcken, a Theodore Wille e ao empréstimo de 1908, que Topik diz
romper cinquenta anos de exclusividade dos Rothschild (p. 72). Testa se a
imprensa lê a pluralidade de credores como ganho de autonomia ou como perda de
controle sobre os estoques.

\item Alinhamento por fração, não por região. Procurar, nas peças sobre a taxa e
sobre o limite de emissão, quem fala por importadores, comércio de praça, bancos
e consumidor urbano, e se esses grupos aparecem do lado da moeda forte, como
Topik afirma (p. 56). Testa se o eixo ortodoxo contra expansionista tem base
social identificável nos jornais.

\item Volta ao ouro ou desvalorização administrada. Procurar se os defensores da
Caixa a apresentam como caminho para a conversibilidade plena ou como instrumento
de taxa baixa para o exportador. Topik sustenta as duas coisas, taxa mais de 50\%
abaixo do par e 40\% do meio circulante conversível em 1912 (pp. 40, 42). Testa
qual justificativa domina e se ela muda entre 1906 e 1910.

\item União contra São Paulo. Procurar acusações de que a Caixa e a valorização
seriam privilégio paulista pago pela União. Topik documenta a objeção da
Associação Comercial da Bahia em 1906 (p. 69) e a versão mineira de Bernardes
(p. 77). Testa se a clivagem federativa já organiza o debate de 1906 no Rio.
\end{dialogo}

\usonocapitulo{Parte II, como principal fonte brasileira sobre enforcement e
sobre a estrutura de interesses por trás da política monetária. Sustenta quatro
afirmações. Que os Rothschild agiram como poder de veto explícito, com ameaça
registrada de perda de soberania em 1898 (p. 36) e oposição declarada a Taubaté
(p. 68), oposição que não era conspiração, já que queriam intervenção estatal
ainda maior (p. 69). Que o monopólio Rothschild se quebra em 1908 e que a
concorrência entre financiadores é parte da história da Caixa e da valorização
(pp. 71-72, 165). Que o Banco do Brasil reorganizado em 1905 é peça de política
cambial subordinada ao ministro da Fazenda, e não banco central independente
(pp. 39-40, 42), o que dá base para tratar a competência da Caixa frente ao
Tesouro e ao Banco do Brasil como disputa institucional real. E que ortodoxia e
intervenção não se opõem no período, com a frase de Bulhões sobre intervir no
mercado de ouro (p. 42) como melhor apoio disponível para o eixo ortodoxo contra
expansionista. Faz ponte com a Parte III pelo mapa de agentes e de frações
(pp. 24, 43, 56). Não serve à Parte IV, porque a Caja argentina não é discutida e
a Caixa não é analisada como \emph{currency board}, nem à Parte I, porque regra e
credibilidade do padrão-ouro não são objeto do livro.}
